\documentclass[10pt,a4paper]{amsart}
\usepackage[margin=19mm]{geometry}
\usepackage{amsmath,amssymb,mathtools}
\usepackage[scr=rsfs,cal=euler]{mathalfa}
\usepackage{xfrac}
\usepackage[unicode,hidelinks,bookmarks=false]{hyperref}
\newcommand{\ie}{i.e.\ }
\newcommand{\cf}{cf.\ }
\newcommand{\resp}{respectively\ }
\newcommand{\Subsection}{\subsection}
\providecommand{\nobreakdash}{\nobreak\hbox{-}}
\setlength{\emergencystretch}{2em}
\multlinegap0pt
\usepackage{bm}
\usepackage{bbm}
\usepackage{dsfont}
\usepackage{xspace}
\usepackage{cancel}
\usepackage{mathtools}
\usepackage{amsthm}
\makeatletter
\@ifundefined{teorema}{\newtheorem{teorema}{Teorema}}{}
\@ifundefined{lema}{\newtheorem{lema}[teorema]{Lemma}}{}
\@ifundefined{proposicio}{\newtheorem{proposicio}[teorema]{Proposition}}{}
\makeatother
\DeclareMathOperator{\Soc}{Soc}
\title[On left braces in which every subbrace is an ideal II]{On left braces in which every subbrace is an ideal II}
\author{A. Ballester-Bolinches, R. Esteban-Romero, L. A. Kurdachenko, V. P\'erez-Calabuig}
\date{Source: arXiv:2601.09580v1 (CC BY 4.0)}
\begin{document}
\maketitle

\noindent\emph{Source and licence.} On left braces in which every subbrace is an ideal II. Version arXiv:2601.09580v1 (CC BY 4.0), distributed under the Creative Commons Attribution 4.0 International licence, as recorded in the arXiv licence metadata for this exact version and shown on the abstract page https://arxiv.org/abs/2601.09580v1. Full rights notice: licenses/arxiv-2601.09580-CC-BY-4.0.txt.\par\smallskip

\noindent\textit{Editorial scope.} Counted unit: the Lema whose source label is \texttt{lema:2.3nou} in the article (source0.tex lines 426--429, source label \texttt{lema:2.3nou}), delivered together with the article's own complete original proof of it (source0.tex lines 431--459, one \texttt{proof} environment of 29 lines). No mathematics is written, repaired or re-derived: statement and proof are the article's own text line for line. Required context retained uncounted: prop:dedekindbraces-Z (lines 221--224); lema:dastd-nozero+Dededkind-><d>=<d>+<dastd> (lines 308--311); lema:infiniteorder->subbridanoideal (lines 379--382); lema:dinSoc2+order-inf+dastd=0->dinS (lines 405--408). Every definition, display and label of those lines is retained unchanged. Omitted from this package and not redistributed: 1--220, 225--307, 312--378, 383--404, 409--425, 430--430, 460--803. Nothing inside a retained excerpt is omitted or altered: every retained body is delivered exactly as the article's lines are written, so no omission receipt of any kind accompanies this package. Citations: the retained excerpts carry no citation command at all, so no bibliography entry of the article is transcribed and no reference is redistributed. Reference adaptation: none was needed and none was made; every reference inside the delivered unit resolves inside it, so no unresolved reference or citation is produced. Printed slips kept literal: no printed word, symbol, hypothesis or number of the counted statement or of its proof was altered. Layout and class: the article's own class is \texttt{article} with packages including fontenc, babel, fancyhdr, inputenc, amsmath, amsthm, amsfonts, amssymb, enumerate, graphicx, xy, url; the wrapper is the lane's pinned \texttt{amsart} editorial wrapper at 10pt on A4, which restates the theorem-like environments the retained text needs and adds the article's own macro definitions that the retained text needs (\texttt{\textbackslash Soc}), with extra wrapper packages (bm, bbm, dsfont, xspace, cancel, mathtools). The delivered numbering is the unit's own. Assets: no retained span contains a figure environment, a graphics-inclusion command, a PDF-page inclusion, a table or a TikZ picture; no image file is copied into this package, the package declares no asset dependency and no image file is redistributed with it. Licence: version arXiv:2601.09580v1 of this article is distributed under the Creative Commons Attribution 4.0 International licence, as recorded in the arXiv licence metadata for this exact version and shown on the abstract page https://arxiv.org/abs/2601.09580v1. A full rights notice is delivered at licenses/arxiv-2601.09580-CC-BY-4.0.txt. Characters: every retained excerpt is pure ASCII, so the delivered engine, XeLaTeX with its default Latin Modern text font, reports no missing character.
\begin{proposicio}
\label{prop:dedekindbraces-Z}
Abelian left braces isomorphic to $\mathbb{Z}$ are the unique Dedekind left braces whose additive group is isomorphic to an infinite cyclic group.
\end{proposicio}

\begin{lema}
\label{lema:dastd-nozero+Dededkind-><d>=<d>+<dastd>}
Suppose that $A$ is a Dedekind left brace and let  $d \in C$ such that $d \ast d \neq 0$. Then, $\langle d \rangle = \langle d \rangle_+ + \langle d\ast d \rangle_+$.  %In particular, if $d$ has finite additive order, $\langle d \rangle$ is finite.
\end{lema}

\begin{lema}
\label{lema:infiniteorder->subbridanoideal}
Let $a\in A$ with infinite additive order. Suppose that $0 \neq c:= a \ast a$ satisfies $c \ast c = a \ast c = c \ast a = 0$. If $c$ has infinite additive order, then $\langle a \rangle = \langle a \rangle_+ \oplus \langle c \rangle_+$ and it includes a subbrace which is not an ideal.
\end{lema}

\begin{lema}
\label{lema:dinSoc2+order-inf+dastd=0->dinS}
Let $A$ be a Dedekind left brace and let $d\in \Soc_2(A)$ with infinite additive order. If $d \ast d = 0$, then $d \in \Soc(A)$.
\end{lema}

\begin{lema}
\label{lema:2.3nou}
Let $A$ be a Dedekind left brace and let $d\in \Soc_2(A)$ with infinite additive order. Then, $d \ast d$ has finite additive order.
\end{lema}

\begin{proof}
Since $d \in \Soc_2(A)$, it follows that $c:= d \ast d \in \Soc(A)$. Suppose that $c$ has infinite additive order. %Thus, we can write $\langle d \rangle = \langle d \rangle_+ \oplus \langle c \rangle_+$. 
Since $\Soc(A)$ is abelian, it follows that $\langle c \rangle = \langle c \rangle_+$ is a subbrace of $A$, and therefore, an ideal. Thus, $(\langle c \rangle_+,\cdot)$ is a normal subgroup of $(A, \cdot)$, so that either $dcd^{-1} = c$ or $dcd^{-1} = -c = c^{-1}$.

Assume $dcd^{-1} = c$. Since $c\in \Soc(A)$, $c \ast c = c \ast d = 0$. Moreover, $d \ast c = dc -d -c = cd -c - d = c\ast d = 0$. Applying, Lemma~\ref{lema:infiniteorder->subbridanoideal}, there exists a subbrace which is not an ideal; a contradiction.

Assume $dcd^{-1} = -c$. Then, it holds that 
\begin{align*}
d \ast c & = dc -d - c = -c + (-c)d -d  = -2c + (-c)d + c - d =\\
& = -2c + (-c)\ast d = -2c.
\end{align*}
On the other hand, since $\Soc_2(A)/\Soc(A)$ is abelian, $2d = d^2s$ for some $s \in \Soc(A)$. Thus, it also holds
\begin{align*}
(2d) \ast (2d) & = (d^2s) \ast (2d) = d^2 \ast \big(s \ast (2d)\big) + s \ast (2d) + d^2 \ast (2d) = \\
& =  d^2 \ast (2d) = 2(d^2 \ast d) = 2\big(d \ast (d \ast d) + 2d \ast d\big) = \\
& = 2(d \ast c + 2c) = 2(-2c + 2c) = 0
\end{align*}
Therefore, by Lemma~\ref{lema:dinSoc2+order-inf+dastd=0->dinS}, $2d\in \Soc(A)$. Hence, $2d + c \in \Soc(A)$ and 
\[ d \ast (2d + c) = 2d \ast d + d \ast c = 0.\]

By Lemma~\ref{lema:dastd-nozero+Dededkind-><d>=<d>+<dastd>}, we can write $\langle d \rangle = \langle d \rangle_+ + \langle c \rangle_+ = \langle d \rangle_+ + \langle 2d + c \rangle_+$. Assume that $\langle d \rangle = \langle d \rangle_+ \oplus \langle 2d + c\rangle_+$. Since $4d+c\in \Soc(A)$, then $\langle 4d +c \rangle = \langle 4d+c \rangle_+$ is an ideal. However, we see that
\[ d \ast (4d + c) = 4(d \ast d) + d \ast c = 4c - 2c = 2c \notin \langle 4d + c \rangle_+ = \langle 2d + 2d+c\rangle_+\]
Hence, we arrive to a contradiction.

Assume that $\langle d \rangle_+ \cap \langle 2d+c\rangle_+ \neq \{0\}$. If $\langle 2d + c \rangle_+ \leq \langle d \rangle_+$, then $\langle d \rangle = \langle d \rangle_+$ is a Dedekind left brace whose additive group is isomorphic to $\mathbb{Z}$. Thus, Proposition~\ref{prop:dedekindbraces-Z} yields that $\langle d \rangle$ is abelian, i.e. $d \ast d = 0$, which is a contradiction. If $\langle d \rangle_+ \leq \langle 2d +c \rangle_+$, then $d \in \Soc(A)$, a contradiction. Therefore, 
\[ \langle d \rangle_+ \neq \langle d \rangle_+ + \langle 2d+c \rangle_+ \neq \langle 2d+c \rangle_+.\] Since $\{0\} \neq \langle d \rangle_+ \cap \langle 2d+c \rangle_+$, the $0$-rank of $\langle d \rangle$ is $1$. Thus, $ \langle d \rangle = \langle v \rangle_+ \oplus \langle u \rangle_+$ where $v$ has infinite additive order and $u$ has finite additive order $n$. Let $l,m \in \mathbb{Z}$ such that  $u = ld + m(2d+c)$. It follows that
\[ d \ast u = d \ast \big(ld+ m(2d+c)\big) = ld \ast d = lc.\]
On the other hand,  $0 = d \ast 0 = d \ast(nu) = n(d\ast u) = nlc$, i.e. $c$ has finite additive order. We arrive to a final contradiction. 
\end{proof}

\end{document}
